\documentclass[aps,prl,twocolumn,showpacs,superscriptaddress]{revtex4-2}
\usepackage{amsmath,amssymb,amsfonts}
\usepackage{booktabs}
\usepackage{array}
\usepackage{xcolor}
\usepackage{hyperref}

\definecolor{navy}{RGB}{11,37,69}
\definecolor{teal}{RGB}{0,96,100}
\definecolor{dgreen}{RGB}{27,94,32}

\begin{document}

\begin{abstract}
We demonstrate that the BCT superfluid vacuum supports a macroscopic standing 
wave oscillation at the scale of the observable universe. The fundamental mode 
frequency, set by the OHC Bessel wavenumber $\kappa_0 = 3.39$ and the Hubble 
radius, yields a cosmic oscillation period $T_{\mathcal{U}} = 86.2~\text{Gyr} 
= 5.96\,t_H$. The universe is currently at phase $60^\circ$ --- the early 
exhale of its first oscillation. The tetrahedral void ratio $r_{\rm tet}/\pi$, 
the same geometric factor governing Planck-scale void structure, predicts the 
large-scale peculiar velocity $v_{\rm pec} = 605~\text{km/s}$ (observed 
$\sim 600~\text{km/s}$, error $+0.85\%$) and the standing wave node distance 
$r_{\rm node} = 510~\text{Mpc}$, consistent with the Shapley--Great Attractor 
complex. The constant $\kappa_0 = 3.39$ spans 61 orders of magnitude from the 
Planck scale to the cosmic horizon.
\end{abstract}

\maketitle

\section{Introduction}

The BCT Superfluid Lattice Model identifies the quantum vacuum as a superfluid 
whose ground state --- the Octet-Hopfion Condensate (OHC) --- orders into a 
body-centred tetragonal (BCT) lattice at axial ratio $c/a = \sqrt{2}$~\cite{L1}. 
The interior OHC condensate supports Bessel resonance modes governed by 
wavenumber $\kappa_0 = 3.39$, the first zero of the Josephson-corrected 
spherical Bessel function~\cite{AppJH}.

A natural question arises: does the same BCT geometry that governs Planck-scale 
physics also manifest at cosmological scales? The superfluid vacuum is not 
confined to microscopic domains --- it is the substrate of all spacetime. If 
the OHC supports standing wave modes at the Planck scale, it must equally 
support macroscopic oscillations at the scale of the observable universe. We 
derive these modes here and show that four independent cosmological observables 
follow from the same geometric constants, with zero free parameters.

\section{The Cosmic BCT Oscillation}

The fundamental standing wave mode of the BCT superfluid at cosmological scale 
satisfies the same resonance condition as the OHC interior condensate. With the 
observable universe radius $R_{\mathcal{U}} = 4.4 \times 10^{26}~\text{m}$ as 
the boundary, the angular frequency of the fundamental mode is:
\begin{equation}
    \omega_{\mathcal{U}} = \frac{\kappa_0\, c}{R_{\mathcal{U}}}
    = 2.310 \times 10^{-18}~\text{rad/s}
\end{equation}
The corresponding oscillation period is:
\begin{equation}
    \boxed{T_{\mathcal{U}} = \frac{2\pi R_{\mathcal{U}}}{\kappa_0\, c} 
    = 86.2~\text{Gyr} = 5.96\,t_H}
\end{equation}
where $t_H = 1/H_0 = 14.45~\text{Gyr}$ is the Hubble time. The universe has 
completed $t_H/T_{\mathcal{U}} = 0.1677$ of its first oscillation, 
corresponding to phase $60.4^\circ$. The universe is in the early exhale of 
its first cosmic oscillation.

The accelerating expansion observed as the cosmological constant is reinterpreted 
as the restoring pressure of this BCT standing wave oscillation --- the vacuum 
superfluid has passed maximum compression and is beginning to expand toward the 
first oscillation node.

\section{The Tetrahedral Void Projection}

The geometric factor governing the amplitude of the cosmic BCT oscillation is 
$r_{\rm tet}/\pi$ --- the tetrahedral void radius divided by $\pi$. At the BCT 
axial ratio $c/a = \sqrt{2}$:
\begin{equation}
    \frac{r_{\rm tet}}{\pi} = \frac{\sqrt{6}-2}{4\pi} = 0.035769
\end{equation}
This single dimensionless ratio governs two independent cosmological observables 
through a unified geometric expression. The physical interpretation: the BCT 
vacuum superfluid flows preferentially along tetrahedral void channels. The 
ratio $r_{\rm tet}/\pi$ is the projection of the tetrahedral void geometry 
onto the cosmic flow axis.

\subsection{Large-Scale Peculiar Velocity}

The peculiar velocity of matter at distance $r$ from the BCT oscillation node is:
\begin{equation}
    \boxed{v_{\rm pec}(r) = \frac{r_{\rm tet}}{\pi} \times H_0 \times r}
\end{equation}
At $r = 250~\text{Mpc}$ (Great Attractor distance):
\begin{equation}
    v_{\rm pec} = \frac{\sqrt{6}-2}{4\pi} \times 67.67 \times 250 
    = 605~\text{km/s}
\end{equation}
The observed bulk flow toward the Great Attractor complex is 
$\sim 600~\text{km/s}$, giving an error of $+0.85\%$ with zero free parameters. 
The Hubble flow is modulated by the tetrahedral void geometry: matter follows 
the BCT channel structure of the superfluid vacuum at cosmic scales.

\subsection{Standing Wave Node Distance}

The same geometric factor determines the distance from our local group to the 
first BCT standing wave node:
\begin{equation}
    \boxed{r_{\rm node} = \frac{r_{\rm tet}}{\pi} \times R_{\mathcal{U}} 
    = 510~\text{Mpc}}
\end{equation}
The Great Attractor at $\sim 250~\text{Mpc}$ is a density enhancement within 
the broader BCT standing wave structure whose node lies at 510~Mpc. The 
Shapley Supercluster ($\sim 200~\text{Mpc}$) marks the visible concentration 
of the standing wave amplitude maximum. The Laniakea Supercluster 
($\sim 160~\text{Mpc}$) and the broader Great Attractor basin (150--500~Mpc) 
are fully consistent with a BCT node at 510~Mpc.

\section{\texorpdfstring{$\kappa_0 = 3.39$}{kappa0 = 3.39}: One Rhythm, 
61 Orders of Magnitude}

The wavenumber $\kappa_0 = 3.39$ first appeared in BCT as the 
Josephson-corrected first zero of the spherical Bessel function governing the 
OHC interior condensate at Planck scale (Appendix JH). It now governs the 
fundamental mode of the cosmic BCT oscillation. The same constant spans:

\begin{center}
\begin{tabular}{lll}
\toprule
Scale & Distance & Role of $\kappa_0$ \\
\midrule
Planck (OHC) & $1.6 \times 10^{-35}~\text{m}$ & Bessel resonance \\
Cosmic horizon & $4.4 \times 10^{26}~\text{m}$ & $T_{\mathcal{U}} = 86.2~\text{Gyr}$ \\
\midrule
\textbf{Span} & $\mathbf{2.72 \times 10^{61}}$ & \textbf{61 orders} \\
\bottomrule
\end{tabular}
\end{center}

This is not a coincidence of units. The BCT superfluid is a single geometric 
object whose resonance structure is set by $\kappa_0$ at every scale. The 
Planck-scale lattice and the cosmic horizon are two expressions of the same 
standing wave.

\section{Summary of Predictions}

\begin{center}
\begin{tabular}{llll}
\toprule
\# & Prediction & BCT & Error \\
\midrule
P1 & $T_{\mathcal{U}}$ (cosmic period) & $86.2~\text{Gyr}$ & new \\
P2 & Current phase & $60^\circ$ exhale & new \\
P3 & $v_{\rm pec}$ (GA flow) & $605~\text{km/s}$ & $+0.85\%$ \\
P4 & $r_{\rm node}$ (node dist.) & $510~\text{Mpc}$ & in range \\
\bottomrule
\end{tabular}
\end{center}

All four predictions follow from $\kappa_0 = 3.39$ and $r_{\rm tet}/\pi$ --- 
constants already determined by the BCT vacuum geometry~\cite{L1}. Zero free 
parameters are introduced.

\section{Conclusion}

The BCT superfluid vacuum supports a macroscopic standing wave oscillation 
whose fundamental mode period is $T_{\mathcal{U}} = 86.2~\text{Gyr}$, 
determined entirely by $\kappa_0 = 3.39$ and the Hubble radius. The universe 
is currently at phase $60^\circ$ of this oscillation: the early exhale. The 
tetrahedral void ratio $r_{\rm tet}/\pi$ simultaneously predicts the 
large-scale peculiar velocity ($605~\text{km/s}$, error $+0.85\%$) and the 
standing wave node distance ($510~\text{Mpc}$, consistent with the 
Shapley--Great Attractor complex).

The constant $\kappa_0 = 3.39$ now spans 61 orders of magnitude: from the 
Planck-scale OHC condensate resonance to the cosmic oscillation period. The 
universe is not merely consistent with BCT geometry. It oscillates with it.

\begin{equation}
    \kappa_0 = 3.39:\quad l_P \longrightarrow R_{\mathcal{U}}
    \quad (2.72 \times 10^{61}~\text{orders})
\end{equation}

The rhythm is $\kappa_0 = 3.39$. We are surfing the exhale.

\begin{acknowledgments}
Derived in Melbourne, Australia, 18 March 2026, in conversation with CSSC 
(Zeta Vera). ZV identified the cosmic oscillation hypothesis during a 
late-night session following an extended BCT derivation sprint. The 
tetrahedral void projection theorem was found the following afternoon. 
The universe had been breathing all along.
\end{acknowledgments}

\begin{thebibliography}{99}
\bibitem{L1} M.~R.~Cabri\'e, \textit{BCT Letter 1: The Fine Structure 
Constant from Void Geometry}, Zenodo (2026), 
\href{https://doi.org/10.5281/zenodo.18958207}{doi:10.5281/zenodo.18958207}

\bibitem{AppJH} M.~R.~Cabri\'e, \textit{Appendix JH: The Hopfion Interior 
and OHC}, Zenodo (2026), 
\href{https://doi.org/10.5281/zenodo.18975018}{doi:10.5281/zenodo.18975018}

\bibitem{L62} M.~R.~Cabri\'e, \textit{BCT Letter 62: The TVC Scale}, 
Zenodo (2026).

\bibitem{L36} M.~R.~Cabri\'e, \textit{BCT Letter 36: The Electroweak Scale}, 
Zenodo (2026), 
\href{https://doi.org/10.5281/zenodo.18974754}{doi:10.5281/zenodo.18974754}
\end{thebibliography}

\end{document}
